% ---
% id: transmitter/rem:da1
% kind: rem
% title: "The one-dimensional retained system"
% summary: >-
%   For a one-dimensional retained system the two source modes impose the
%   same condition and the licence excludes no attack. A trusted non-empty
%   source segment there for which no certificate exists rests on the
%   supplier obligation A-LINK, a case no artefact reaches.
% prominence: standard
% status: accepted
% uses: [def:source-spec, def:srcert]
% mentions: [ass:src, rem:source-mode, req:licence-scope, def:source-segment, def:source-link, thm:certificate-sufficiency, lem:ensemble-presentation, reductions/def:chain, reductions/thm:two-sided, reductions/lem:source-transfer, protocol-compiler/ass:link, protocol-compiler/def:chaincond, protocol-compiler/thm:chain-consumption, meat-core/rem:modes-da1, transmitter/decoy/def:tagged-source]
% assumes: []
% refs: []
% provenance: newly-derived
% lean:
% ---

\begin{remark}
\label{rem:da1}
Suppose $d_A=1$. Then $\Seq{A}=\{[1]\}$, so the only marginal a
trusted-mode specification can carry is $\tau_A=[1]$; and for any state
$\omega$ on $A$ and the \tref[virtual transmitted system]{transmitted system} the $A$-marginal $\omega_A$ is
the $1\times1$ matrix $\tr\omega$, so the conditions $\omega_A=\tau_A$ and
$\tr\omega=1$ are the same condition on the same number. The two source
modes, \tref{trusted-source mode} and \tref{untrusted-source mode}, therefore impose the same
condition on every consumer, and what the source-replacement certificate
licenses at $d_A=1$ is the condition $\tr\omega=1$, which every normalized
state satisfies: it excludes nothing, and Assumption~\ref{ass:src} is
\emph{vacuous} --- it deletes no attack. That is why the certificate of
Definition~\ref{def:srcert} is demanded only at $d_A\ge2$
(Requirement~\ref{req:licence-scope}). A specification at $d_A=1$ is also
the shape in which a protocol with no trusted source is declared: the
consuming declaration folds Alice's measured system into the receiver-side
systems and names, in its partition of the outcome among the parties, which
components are hers (\xref{meat-core/rem:modes-da1}). That shape is the
\emph{safe} direction (Remark~\ref{rem:source-mode}): it removes an assumption
rather than adding one.

\emph{A non-empty source segment at $d_A=1$.} Every such specification that
an artefact of this module can carry is covered by theorems, and the one case
left to an assumption is reached by no artefact. In detail, the theorem route
applies to a non-empty segment whose consumer declares the ensemble form
(Definition~\ref{def:srcert}, ``What a consumer reads of a presentation'')
and which passes every condition that admits it: the admission rule on the
kind, (S1) to (S4) and (A1) to (A3) of Definition~\ref{def:source-segment},
the refusals of Definition~\ref{def:source-link}\,(d) and (f), and the
admission rule of \xref{reductions/def:chain}, which admits the segment only
in trusted-source mode and only where every entry of the constraint list is
licensed, as \xref{protocol-compiler/def:chaincond} decides it. There the
security statement is carried across the segment by the source identity
\eqref{eq:srcmap} at the descriptors of the emitted links, by
Theorem~\ref{thm:certificate-sufficiency}\,(iii) with (SR1) to (SR3) for the source-replacement
recasting, and by \xref{reductions/thm:two-sided} for the composition with
the receiver segment; at $d_A\ge2$ Assumption~\ref{ass:src} supplies the
attestation. The degenerate declaration of the decoy submodule
(Definition~\ref{def:tagged-source}) is covered by the same route: with one
polarisation, intensity set $\{0\}$, its \tref{entropy cutoff} and
\tref{source cutoff} both zero and
all weight on the zero-photon sector it has one one-dimensional sector, hence
$d_A=1$; its ensemble presentation supplies (SR1) and (SR2) as theorems
(Lemma~\ref{lem:ensemble-presentation}\,(i)), its defined source state and
marginal are both $[1]$, and (SR3) is vacuous, a trace-preserving map being
unable to act nontrivially on a one-dimensional retained register. The
hypotheses of \xref{reductions/lem:source-transfer}\,(ii), (iii) and
\xref{reductions/thm:two-sided} therefore hold there as well.

\emph{The case left to an assumption.} A trusted specification at $d_A=1$
that declares no certificate, with a non-empty segment and a virtual
ensemble that is not constant over the letters with $p_x>0$, has none to
declare: with a one-dimensional retained register (SR2) would force all
positively weighted virtual states to coincide. The theorem route does not
cover that case. Its transfer across the segment is the source clause of the
protocol compiler's supplier obligation A-LINK
(\xref{protocol-compiler/ass:link}), declared on that module's
chain-condition definition, \xref{protocol-compiler/def:chaincond}, and on
\xref{protocol-compiler/thm:chain-consumption}\,(v$'$), never discharged, and emitted in
its assumption ledger. No artefact reaches the case: (A3) of
Definition~\ref{def:source-segment} accepts a segment only inside an artefact
of the admitting module, and the one admitted constructor kind emits an
ensemble presentation, which supplies the certificate described above.
\end{remark}
